\documentclass[12pt]{article}

\setlength{\parindent}{0pt}
\setcounter{secnumdepth}{3}
\usepackage[margin=1in]{geometry}

% packages
\usepackage{amsmath, amsfonts, amsthm, amssymb}
\usepackage{graphicx, float}
\usepackage{mathtools}
\usepackage{titlesec}
\usepackage{interval}
\usepackage{hyperref}
\usepackage{siunitx}
\usepackage{titling}
\usepackage{vwcol}
\usepackage{setspace}
\usepackage{empheq}
\usepackage{cancel}
\usepackage{esdiff}
\usepackage{multicol}
\usepackage{mdframed}
\usepackage{esdiff}
\usepackage{tikzsymbols}
\usepackage{multicol}
\usepackage{tikz}
\usepackage{varwidth}
\usepackage{pgfplots}
\usepackage{fancyhdr}
\usepackage{enumitem}
\usepackage{tcolorbox}

\intervalconfig {
	soft open fences
}

% header
\pagestyle{fancy}
\fancyhf{}
\lhead{Ethan Fan}
\rhead{PCH}
\cfoot{\thepage}

% section format
\titleformat{\section}
  {\large \bfseries}
  {}
  {0em}
  {}
\titlespacing*{\section}{0pt}{0.5em}{0.3em}

\titleformat{\subsection}[runin]
  {\bfseries}
  {}
  {0em}
  {}
\titlespacing*{\subsection}{0pt}{0pt}{0em}

\titleformat{\subsubsection}[runin]
  {\itshape}
  {(\roman{subsubsection})}
  {0em}
  {}

% commands
\newcommand{\alignedintertext}[1]{%
  \noalign{%
    \vskip\belowdisplayshortskip
    \vtop{\hsize=\linewidth#1\par
    \expandafter}%
    \expandafter\prevdepth\the\prevdepth
  }%
}

\newcommand*{\paren}[1]{\ensuremath\left(#1\right)}
\newcommand*{\limit}[2][x]{\ensuremath{\displaystyle\lim_{#1\to#2}}}
\newcommand*{\Limit}[3][x]{\ensuremath{\displaystyle\lim_{#1\to#2}\left[#3\right]}}
\newcommand*{\deriv}[1][x]{\ensuremath{\dfrac{\mathrm{d}}{\mathrm{d}#1}}}
\newcommand*{\Deriv}[2][x]{\ensuremath{\dfrac{\mathrm{d}}{\mathrm{d}#1}\left[#2\right]}}
\newcommand{\dydx}{\dfrac{dy}{dx}}
\newcommand{\dydt}{\dfrac{dy}{dt}}
\newcommand{\dxdt}{\dfrac{dx}{dt}}
\newcommand*{\iinteg}[2][x]{\ensuremath{\displaystyle\int #2\;\mathrm{d}#1}}
\newcommand*{\dinteg}[4][x]{\ensuremath{\displaystyle\int_{#2}^{#3}#4\;\mathrm{d}#1}}
\newcommand*{\abs}[1]{\ensuremath{\left|#1\right|}}
\newcommand*{\eps}{\varepsilon}
\newcommand*{\real}{\ensuremath\mathbb{R}}
\newcommand*{\integer}{\ensuremath\mathbb{Z}}
\newcommand*{\posint}{\ensuremath\mathbb{Z^+}}
\newcommand*{\cbrt}[1]{\ensuremath{\sqrt[3]{#1}}}
\newcommand*{\lheqzero}{\ensuremath{\underset{\text{L'H}}{\overset{\left[\frac00\right]}{=}}}}
\newcommand*{\lheqinfty}{\ensuremath{\underset{\text{L'H}}{\overset{\left[\frac{\infty}{\infty}\right]}{=}}}}
\newcommand{\tor}{\text{ or }}
\newcommand{\floor}[1]{\lfloor #1 \rfloor}

\newcommand{\vem}{\vspace{0.5em}}
\newcommand{\example}[1]{\section{Example #1}}
\newcommand{\problem}[1]{\section{Problem #1}}
\newcommand{\subproblem}[1]{\subsection{(#1)} \,}

\DeclareMathOperator{\dne}{DNE}
\DeclareMathOperator{\sgn}{sgn}

\newcommand{\definition}[1]{\begin{tcolorbox}[colback=red!5!white,colframe=red!75!black,parbox=false] #1 \end{tcolorbox}}

\newcommand{\theorem}[2]{\begin{tcolorbox}[title={#1},colback=blue!5!white,colframe=blue!75!black,parbox=false] #2 \end{tcolorbox}}

\allowdisplaybreaks
% \postdisplaypenalty=100000

% document
\begin{document}

\vspace*{0cm}

% opening
\begin{center}
    {\Large \bfseries Problem Set 76} \\[0.5cm]
    {\large Optimization Part 1} \\[0.35cm]
    {\normalsize May 15, 2026}
\end{center}

% problems

\problem{1}
Find two positive numbers whose product is 100 and whose sum is minimum.
\begin{gather*}
    \begin{cases}
        S=x+y\\
        xy=100
    \end{cases}\\
    S=x+\frac{100}{x}, \quad x>0\\
    S'=1-\frac{100}{x^2}\\
    S'=0 \implies x=10\\
    \boxed{x=y=10}
\end{gather*}
The two positive integers are 10 and 10.

\problem{2}
Find the dimensions of a rectangle with perimeter 100 m whose area is as large as possible.
\begin{gather*}
    \begin{cases}
        A=xy\\
        2x+2y=100
    \end{cases}\\
    A=x(50-x)=-x^2+50x, \quad x>0\\
    A'=-2x+50\\
    A'=0 \implies \boxed{x=y=25}
\end{gather*}
The dimensions of the rectangle is 25 m by 25 m.

\problem{3}
A cylindrical can is to be made to hold 1 L of oil. Find the dimensions that will minimize the cost of the metal to manufacture the can.
\begin{gather*}
    \begin{cases}
        \pi r^2h=1\\
        A=2\pi r^2+2\pi rh
    \end{cases}\\
    A=2\pi r^2+\frac{2}{r}, \quad r>0\\
    A'=4\pi r-\frac{2}{r^2}\\
    A'=0 \implies \boxed{r=\cbrt{\frac{1}{2\pi}}, h=2\cbrt{\frac{1}{2\pi}}}
\end{gather*}
The cylindrical can has a radius of $r=\cbrt{\dfrac{1} {2\pi}}$ dm and a height of $h=2\cbrt{\dfrac{1}{2\pi}}$ dm.

\problem{4}
 A model used for the yield $Y$ of an agricultural crop as a function of the nitrogen level $N$ in the soil
(measured in appropriate units) is
\[
Y=\frac{kN}{1+N^2}
\]
where $k$ is a positive constant. What nitrogen level gives the best yield?
\begin{gather*}
    Y=\frac{kN}{1+N^2}, \quad N>0\\
    Y'=\frac{k(1+N^2)-2N\cdot kN}{(1+N^2)^2}=\frac{k(1-N^2)}{(1+N^2)^2}\\
    Y'=0 \implies \boxed{N=1}
\end{gather*}
The optimal nitrogen level $N$ is 1.

\problem{5}
 A farmer wants to fence an area of 1.5 million square feet in a rectangular field and then divide it in half
with a fence parallel to one of the sides of the rectangle. How can he do this so as to minimize the cost
of the fence?
\begin{gather*}
    \begin{cases}
        A=lh\\
        P=3l+2h
    \end{cases}\\
    P=3l+\frac{2A}{l}, \quad l>0\\
    P'=3-\frac{2A}{l^2}\\
    P'=0 \implies l=\sqrt{\frac{2A}{3}} \implies \boxed{l=1000, h=1500}
\end{gather*}
To minimize the cost, the dimensions of the fence are 1000 ft by 1500 ft.

\problem{6}
Find the point on the line $y = 4x + 7$ that is closest to the origin.
\begin{gather*}
    \begin{cases}
        y=4x+7\\
        d=\sqrt{x^2+y^2}
    \end{cases}\\
    d=\sqrt{x^2+16x^2+56x+49}\\
    d=\sqrt{17x^2+56x+49},\quad d>0\\
    d'=0\implies 34x+56=0 \implies \boxed{x=-\frac{28}{17}}
\end{gather*}
The point is $\paren{-\dfrac{28}{17},\dfrac{7}{17}}$.

\problem{7}
Find the points on the ellipse $4x^2 + y^2 = 4$ that are farthest away from the point $(1, 0)$.
\begin{gather*}
    \begin{cases}
        4x^2 + y^2 = 4\\
        d^2=(x-1)^2+y^2
    \end{cases}\\
    d^2=(x-1)^2+4-4x^2\\
    d^2=-3x^2-2x+5\\
    (d^2)'=-6x-2\\
    (d^2)'=0 \implies \boxed{x=-\frac{1}{3}}
\end{gather*}
The points are $\paren{-\dfrac{1}{3},\dfrac{4\sqrt{2}}{3}}$ and $\paren{-\dfrac{1}{3},-\dfrac{4\sqrt{2}}{3}}$.

\problem{8}
Find the dimensions of the rectangle of the largest area that can be inscribed in a circle of radius $r$.
\begin{gather*}
    \begin{cases}
        x^2+y^2=r^2\\
        A=xy
    \end{cases}\\
    A=x\sqrt{r^2-x^2}\\
    A'=\sqrt{r^2-x^2}-\frac{x^2}{\sqrt{r^2-x^2}}\\
    A'=0 \implies r^2=2x^2 \implies \boxed{x=\frac{r}{\sqrt{2}}}
\end{gather*}
The dimensions of the rectangle are $\sqrt{2}r$ by $\sqrt{2}r$.

\problem{9}
Find the dimensions of the rectangle of largest area that can be inscribed in an equilateral triangle of
side $L$ if one side of the rectangle lies on the base of the triangle.
\begin{gather*}
    \begin{cases}
        A=xy\\
        y=-\sqrt{3}x+\dfrac{\sqrt{3}L}{2}
    \end{cases}\\
    A=x(-\sqrt{3}x+\frac{\sqrt{3}L}{2})=-\sqrt{3}x^2+\frac{\sqrt{3}Lx}{2}\\
    A'=-2\sqrt{3}x+\frac{\sqrt{3}L}{2}\\
    A'=0 \implies \boxed{x=\frac{L}{4},y=\frac{L\sqrt{3}}{4}}
\end{gather*}
The dimensions of the rectangle are $\dfrac{L}{2}$ by $\dfrac{\sqrt{3}L}{4}$.

\problem{10}
Find the dimensions of the isosceles triangle of largest area that can be inscribed in a circle of radius $r$.
\begin{gather*}
    \begin{cases}
        A=\dfrac{1}{2}x^2 \sin 2\theta\\
        x=2r \cos \theta
    \end{cases}\\
    A=\frac{1}{2}(2r\cos \theta )^2\sin 2\theta=2r^2 \cos^2 \theta \sin 2\theta, \quad \theta >0\\
    A'=2r^2(2\cos \theta (- \sin \theta)\sin 2\theta + 2\cos^2 \theta \cos 2\theta )\\
    A'=4r^2\cos \theta \cos 3\theta \\
    A'=0 \implies \boxed{\theta =\frac{\pi}{6}}
\end{gather*}
The triangle is equilateral with a side length of $\sqrt{3}r$.

\problem{11}
 A right circular cylinder is inscribed in a sphere of radius $r$. Find the largest possible volume of such a
cylinder.
\begin{gather*}
    \begin{cases}
        v=2\pi x^2y\\
        x^2+y^2=r^2
    \end{cases}\\
    v=2\pi y(r^2-y^2) ,\quad 0<y<r\\
   v'=2\pi r^2-6\pi y^2\\
   v'=0 \implies \boxed{y=\frac{r}{\sqrt{3}},x=r\sqrt{\frac{2}{3}}}
\end{gather*}
The volume of the cylinder is $\dfrac{4\pi r^3}{3\sqrt{3}}$.

\problem{12}
A right circular cylinder is inscribed in a sphere of radius $r$. Find the largest possible surface area of
such a cylinder.
\begin{gather*}
    \begin{cases}
        A=2\pi x^2+4\pi xy\\
        x^2+y^2=r^2
    \end{cases}\\
    A=2\pi x^2+4\pi x\sqrt{r^2-x^2},\quad 0<x<r\\
    A'=4\pi x + 4\pi (\sqrt{r^2-x^2}-\frac{x^2}{\sqrt{r^2-x^2}})\\
    A'=0 \implies x=\frac{r^2-2x^2}{\sqrt{r^2-x^2}}\\
    x^2(r^2-x^2)=4x^4-4r^2x^2+r^4\\
    5x^4-5r^2x^2+r^4=0\\
    \boxed{x=r\sqrt{\frac{5+\sqrt{5}}{10}}, y=r\sqrt{\frac{5-\sqrt{5}}{10}}}\\
    A=2\pi x^2 + 4\pi xy\\
    A= 2 \pi r^2 \cdot \frac{5+\sqrt{5}}{10} + 4\pi \cdot \frac{r^2}{\sqrt{5}}\\
    \boxed{A=\pi r^2 (1+\sqrt{5})}
\end{gather*}
The largest possible surface area of the cylinder is $\pi r^2 (1+\sqrt{5})$.

\problem{13}
 A piece of wire 10 m long is cut into two pieces. One piece is bent into a square and the other is bent
into an equilateral triangle. How should the wire be cut so that the total area enclosed is \vem

\subproblem{a}a maximum?
\begin{gather*}
    x=0 \implies A=\frac{\sqrt{3}}{36}\cdot 10^2\\
    x=10 \implies A=\frac{1}{16}\cdot 10^2\\
    \boxed{x=10}
\end{gather*}
The maximum area is attained when the entire wire is bent into a square. \vem

\subproblem{b}a minimum?
\begin{gather*}
    \begin{cases}
        10=x+y\\
        A=(\dfrac{x}{4})^2+\dfrac{\sqrt{3}}{4}\cdot (\dfrac{y}{3})^2
    \end{cases}\\
    A=\frac{x^2}{16}+\frac{\sqrt{3}}{36}\cdot (10-x)^2,  \quad 0\le x\le 10\\
    A'=\frac{x}{8}-\frac{\sqrt{3}}{18}\cdot (10-x)\\
    A'=0 \implies x=\frac{80\sqrt{3}}{18+8\sqrt{3}}\\
    \boxed{x=\frac{120\sqrt{3}-160}{11}}
\end{gather*}
The minimum area is attained when a segment of length $\dfrac{120\sqrt{3}-160}{11}$ m is cut and bent into a square, while the rest is bent into a triangle.

\problem{14}
A cone-shaped drinking cup is made from a circular piece of paper of radius $R$ by cutting out a sector
and joining the edges $CA$ and $CB$. Find the maximum capacity of such a cup.
\begin{gather*}
    \begin{cases}
        V=\dfrac{1}{3} \pi r^2h\\
        r^2+h^2=R^2
    \end{cases}\\
    V=\frac{\pi}{3}(R^2h-h^3),\quad 0<h<R\\
    V'=\frac{\pi}{3}(R^2-3h^2)\\
    V'=0 \implies h=\frac{R}{\sqrt{3}}, r=\frac{\sqrt{2}R}{\sqrt{3}}\\
    V=\frac{\pi}{3}(\frac{2R^2}{3})(\frac{R}{\sqrt{3}})=\boxed{\dfrac{2\pi\sqrt{3}R^3}{27}}
\end{gather*}
The maximum capacity of the cup is $\dfrac{2\pi\sqrt{3}R^3}{27}$.

\problem{15}
A cone with height $h$ is inscribed in a larger cone with height $H$ so that its vertex is at the center of the base of the larger cone. Show that the inner cone has maximum volume when $h = \dfrac{1}{3}H$.
\begin{gather*}
    \begin{cases}
        V=\dfrac{1}{3}\pi r^2h\\
        r=\dfrac{R(H-h)}{H}
    \end{cases}\\
    V=\frac{\pi R^2(h^3-2Hh^2+H^2h)}{3H^2}, \quad 0<h<H\\
    V'=\frac{\pi R^2}{3H^2}(3h^2-4Hh+H^2)\\
    V'=0 \implies \boxed{h=\frac{H}{3}}
\end{gather*}
The inner cone has maximum volume when $h=\dfrac{H}{3}$.

\problem{16}
If a resistor of $R$ ohms is connected across a battery of $E$ volts with internal resistance $r$ ohms, then the
power (in watts) in the external resistor is
\[
P=\frac{E^2R}{(R+r)^2}
\]
If $E$ and $r$ are fixed but $R$ varies, what is the maximum value of the power?
\begin{gather*}
    P'=\frac{E^2(R+r)^2-2E^2R(R+r)}{(R+r)^4}\\
    P'=0 \implies R=r\\
    \boxed{P=\frac{E^2}{4r}}
\end{gather*}
The maximum value of the power is $\dfrac{E^2}{4r}$ watts.

\problem{17}
In a beehive, each cell is a regular hexagonal prism, open at one end with a trihedral angle at the other end as in the Figure 2. It is believed that bees form their cells in such a way as to minimize the surface area for a given volume, thus using the least amount of wax in cell construction. Examination of these cells has shown that the measure of the apex angle $\theta$ is amazingly consistent. Based on the geometry of the cell, it can be shown that the surface area is given by
\[
S = 6sh - \frac{3}{2}s^2 \cot \theta + 3s^2 \frac{\sqrt{3}}{2} \csc \theta
\]
where $s$ the length of the sides of the hexagon, and $h$, the height, are constants. \vem

\subproblem{a} Calculate $\dfrac{dS}{d\theta}$.
\[
\boxed{\dfrac{dS}{d\theta}=\frac{3}{2}s^2 \csc^2\theta - 3s^2 \frac{\sqrt{3}}{2} \csc \theta \cot \theta }
\]

\subproblem{b} What angle should the bees prefer?
\begin{gather*}
    \dfrac{dS}{d\theta}=0 \implies \frac{3}{2}s^2\csc \theta (\csc \theta -\sqrt{3}\cot \theta )=0\\
    \csc \theta =\sqrt{3}  \cot \theta \\
    \cos \theta =\frac{\sqrt{3}}{3}\\
    \boxed{\theta =\arccos \frac{\sqrt{3}}{3}}
\end{gather*}
The preferred angle is $\arccos \dfrac{\sqrt{3}}{3}$.

\subproblem{c} Determine the minimum surface area of the cell (in terms of $s$ and $h$).\vem

\textit{Note: Actual measurements of the angle $\theta$ in beehives have been made, and the measures of this angles seldom differ from the calculated value by more than 2 deg.}
\begin{align*}
    S &= 6sh - \frac{3}{2}s^2 \cot \theta + 3s^2 \frac{\sqrt{3}}{2} \csc \theta\\
    &= 6sh-\frac{3}{2}s^2\cdot \frac{1}{\sqrt{2}}+3s^2\frac{\sqrt{3}}{2}\cdot \frac{\sqrt{3}}{\sqrt{2}}\\
    &=\boxed{6sh+\frac{3\sqrt{2}}{2}s^2}
\end{align*}
The minimum surface area of the cell is $6sh+\dfrac{3\sqrt{2}}{2}s^2$.

\problem{18}
An oil refinery is located on the north bank of a straight river that is 2 km wide. A pipeline is to
be constructed from the refinery to storage tanks located on the south bank of the river 6 km east of
the refinery. The cost of laying pipe is \$$400, 000$/km over land to a point $P$ on the north bank and
\$$800, 000$/km under the river to the tanks. To minimize the cost of the pipeline, where should $P$ be
located?
\begin{gather*}
    C=800000\sqrt{4+x^2}+400000(6-x), \quad 0\le x\le 6\\
    C'=\frac{800000x}{\sqrt{4+x^2}}-400000\\
    C'=0 \implies 2x=\sqrt{4+x^2}\\
    4x^2=4+x^2\\
    \boxed{x=\frac{2}{\sqrt{3}} $ km$}
\end{gather*}
Point $P$ is located $6-\dfrac{2}{\sqrt{3}}$ km west of the storage tanks. 

\problem{19}
The illumination of an object by a light source is directly proportional to the strength of the source and
inversely proportional to the square of the distance from the source. If two light sources, one three times
as strong as the other, are placed 10 ft apart, where should an object be placed on the line between the
sources so as to receive the least illumination?
\begin{gather*}
    E = \frac{I}{d^2}+\frac{3I}{(10-d)^2},\quad 0<d<10\\
    E'=-\frac{2I}{d^3}+\frac{6I}{(10-d)^3}\\
    E'=0 \implies 3d^3=(10-d)^3\\
    \cbrt{3}d=10-d\\
    \boxed{d=\frac{10}{\cbrt{3}+1}}
\end{gather*}
The object is placed $\dfrac{10}{\cbrt{3}+1}$ ft away from the weaker light source. 

\problem{20}
Show that of all the isosceles triangles with a given perimeter, the one with the greatest area is equilateral.
\begin{gather*}
    \begin{cases}
        P=2a+b\\
        A=\sqrt{s(s-a)^2(s-b)}
    \end{cases}\\
    A=\sqrt{s(s-a)^2(2a-s)}\\
    f(a)=(a^2-2as+s^2)(2a-s), \quad 0<a<s\\
    f(a)=2a^3-5sa^2+4s^2a-s^3\\
    f'(a)=6a^2-10sa+4s^2\\
    f'(a)=0 \implies a=\frac{2s}{3}\\
    \boxed{a=b=\frac{P}{3}}
\end{gather*}
Hence, the triangle with the greatest area is equilateral.

\end{document}

